\chapter{The afterword}
\label{ch:afterword}

\begin{aside}
Some final thoughts after all those pages.
\end{aside}

\section{Why this book}

The terminal is small. The code in this book is mostly
small. But the ideas are not. Every decision in \maya{}
trades off something against something else; this book
tried to make those trade-offs visible.

\section{What was hard to write}

Layout was hardest --- it's mostly intuition for users
and mostly arithmetic for implementers; bridging that gap
in words took attempts. Animations were second hardest;
they feel obvious but are subtle once you implement.

\section{What was easy}

Patterns came naturally. Once you've built a few apps in
\maya{}, the patterns suggest themselves.

\section{What was missing}

A real treatment of BiDi. Better coverage of mobile
terminals. More material on testing async flows. Maybe
later editions.

\section{Who helped}

Open-source contributors. The terminal protocol authors
across decades. Other library authors whose work
\maya{} learned from. The user who reported bug \#47.

\section{What's next for the book}

It will follow \maya{}. New chapters as new features land.
Errata as bugs are found. Patches welcome.

\section{Recap}

\begin{recap}
\begin{itemize}[leftmargin=*]
  \item Small medium, big ideas.
  \item Trade-offs made visible.
  \item Living document; patches welcome.
\end{itemize}
\end{recap}

\medskip

\noindent Until next time.
